\section{Relation to the Note}


This section organizes the relation between this paper and the companion note series.

This paper is not a replacement for the note series.

The note series is an intuitive reading on consciousness, beginning from experience, metaphor, unease, and everyday sensation.

This paper, by contrast, is the theoretical foundation for fixing the structure behind that.

The role of this paper is therefore not to overwrite the note series with theory but to reposition what was intuitively shown in the note series within the structure of temporal non-closure, speed-scale difference, log referencing, and meaning gradient.

\subsection{The Note as Intuitive Layer}


In the note series, what readers should first encounter is not equations.

What readers should first encounter is a sense of unease about their own conscious experience.

Why does the body begin moving before one thinks one has decided? Why does a sensation change slightly the moment it is put into words? Why do strong memories and emotions repeatedly return to consciousness? Why does consciousness feel as though it is inside oneself, yet cannot be found when searched for?

The note series handles these questions intuitively.

Theoretical terms can be minimal there.

What matters is that readers feel: "Perhaps consciousness is less something generated and more something that opens."

\subsection{The Present Paper as Structural Layer}


This paper, on the other hand, structures the intuitions presented in the note series.

In this paper, consciousness has been formalized as follows:

\begin{itemize}
\item Consciousness is not a product but a log-referencing state.
\item Consciousness is not a single site but a temporal non-closure regime.
\item Consciousness depends on speed difference $\Delta v$ between fast and slow layers.
\item Consciousness requires coupling viscosity $\mu$ and internal delay $\tau$.
\item Log referencing is directed by meaning gradient $\nabla\Phi$.
\item Effective torque-like quantity $\mathcal{T}_{\Phi}$ represents the coupling of speed difference and meaning gradient.
\item Qualia are not the substance of consciousness but the live foreground at the sedimentation boundary.
\item Measurement converts that live boundary into residue.
\item Consciousness is non-local; it opens when multiple local substrates enter a specific temporal relation.
\end{itemize}

These need not be brought to the foreground in the note series.

But they must be retained as the structure behind the note series.

\subsection{What Should Flow Back to the Note}


What should flow back from this paper to the note series is not all the equations.

What should flow back is the minimal structure that helps readers understand.

Concretely:

\begin{enumerate}
\item The expression that consciousness "opens" rather than "is produced."
\item The intuition that consciousness opens between fast processing and slow memory sedimentation.
\item The organization of the ego not as subject but as transformation point of speed difference.
\item The organization of qualia not as stored objects but as live foreground at the sedimentation boundary.
\item The organization that measurement and introspection do not capture qualia directly but log them.
\item The organization that consciousness obtains not as a single brain point but as a temporal condition.
\end{enumerate}

These are sufficiently intuitive for note series readers while also connecting to the theoretical foundation of this paper.

\subsection{What Should Not Flow Back Directly}


Conversely, some things should not be brought directly from this paper into the note series.

First, equations should not be brought in excessively. $\Delta v$, $\mu$, $\tau$, $\mathcal{T}_{\Phi}$, $\nabla\Phi$ are important structural variables in this paper, but foregrounding all of them in the note series breaks the flow as a reading.

Second, pathological deviations should not be treated as diagnoses. The pathologies of this paper are directions of deviation from the consciousness window, not psychiatric classifications. Even when treating these in the note series, they must be handled carefully as deformations of conscious states, not medical explanations.

Third, IFGT/VED terminology should not be brought too prominently to the foreground. The role of the note series is not to make readers understand theory names but to let readers experience structure. Theory should support from behind, not advance to the front.

\subsection{Layered Translation}


The relation between this paper and the note series is not simple summarization.

It is translation across layers.

At L3, consciousness is written:

\begin{equation}
\text{Consciousness} \sim \text{Log Referencing}\!\left(\Delta v, \mu, \tau, \mathcal{T}_{\Phi}, \nabla\Phi\right)
\end{equation}

But in the note series, there is no need to present this directly.

In the note series, it can be restated as:

Consciousness is the passage that temporarily opens between fast-flowing processing and slowly sinking memory.

This translation is important. The same structure is shown at different resolutions.

This is the relation between the note series and this paper.

\subsection{Relation to Appendices}


Between the note series and this paper, supplementary materials and intermediate explanatory layers are placed.

This paper fixes structure as L3.

Supplementary materials explain that structure to readers as L2.

Where necessary, L2.5 serves as a bridge that returns the meaning of equations to intuition.

The note series is the experiential layer readers encounter first, as L1.

Not conflating these layers is important.

L3 must not dominate L1.

L1 must not blur L3.

Each layer has a different role.

The purpose of this paper is not to harden the text of L1 but to fix the structure behind it so that L1 can write with freedom and ease.

\subsection{Summary}


The relation between this paper and the note series is repositioning, not replacement.

\begin{itemize}
\item The note series is the layer of experience and intuition.
\item This paper is the layer of structure and equations.
\item Not everything in this paper needs to be brought into the note series.
\item What should flow back to the note series is the minimal structure that helps readers' intuition.
\item Theoretical terms and equations support through supplementary materials and intermediate layers when needed.
\item L3 does not overwrite L1.
\item L3 supports L1 from behind.
\end{itemize}

This paper is therefore not the completed version of the note series.

This paper is the skeleton behind the note series — enabling the note series to breathe freely as a reading.
